\clearpage
\SourceLeaf{183}
Ähnliches gilt für die Karte der Gesamtintensität, die Herr Dr. \textsc{Goldschmidt} nach den Tabellen erstellt hat. Dabei haben sich zwei Maximalpunkte und ein Kreuzungspunkt in der nördlichen Hemisphäre, ein Maximalpunkt in der südlichen Hemisphäre sowie zwei Minimalpunkte und zwei Kreuzungspunkte in der mittleren Zone ergeben.

An ähnlichen, auf der Theorie beruhenden Karten für die Inklination, die horizontale Intensität, die drei Komponenten der erdmagnetischen Kraft und die Verteilung der magnetischen Flüssigkeiten auf der Erdoberfläche. die stellvertretend für die tatsächliche Verteilung im Inneren gelten kann [siehe Artikel 32], wird bereits gearbeitet. Wir hoffen, sie dem nächsten Band der \textit{Resultate} beifügen zu können.

[Alle hier genannten Karten sowie Tabellen mit den von $5$ zu $5^0$ Breite und von $10$ zu $10^0$ Länge berechneten Werten sowohl der in Artikel 27 mit $\dfrac VR$, $X$, $Y$, $Z$ bezeichneten Größen als auch der Deklination, der Inklination, der Gesamtintensität und der horizontalen Intensität wurden unter dem Titel ‘\textit{Atlas des Erdmagnetismus nach den Elementen der Theorie entworfen}’ als Ergänzung zu den Resultaten aus den Beobachtungen des magnetischen Vereins unter Mitwirkung von \textsc{C. W. B. Goldschmidt} von \textsc{Carl Friedrich Gauss} und \textsc{Wilhelm Weber} in Leipzig 1840 herausgegeben. Da von diesem Atlas zurzeit noch genügend Exemplare vorhanden sind, ist dem vorliegenden Wiederabdruck der allgemeinen Theorie des Erdmagnetismus nur die Karte für die Werte von $\dfrac VR$ beigefügt.]

[Zusätzlich zu den in der Abhandlung angegebenen Vergleichen der Formeln mit den Beobachtungen an 99 Orten enthalten die \textit{Resultate} und der genannte \textit{Atlas} die hier zusammengestellten Vergleiche für 44 weitere Orte.

Die Angaben für die Insel Zafarine, für Toulon und für den Ort auf $70^0\;53'$ nördlicher Breite und $170^0$ Länge finden sich im Atlas.

Die übrigen Vergleiche wurden von \textsc{B. Goldschmidt} berechnet. Er hat zu den Beobachtungen die unten folgenden Nachweise in den \textit{Resultaten} für 1840 und 1841 angegeben.]
\clearpage
\SourceLeaf{184}
\begingroup
\fontsize{9}{10.2}\selectfont
\setlength{\tabcolsep}{1.5pt}
\renewcommand{\arraystretch}{1.02}
\begin{center}
\begin{tabular}{rl|r|r||r|r|r}
\hline
&&Breite&Länge&\multicolumn{3}{c}{Deklination}\\
&&&&Berechn.&Beobacht.&Differenz\\\hline
$1$&Auf dem Eise&$+70^0\;53'$&$170^0\;0'$&$-16^0\;47'$&$-18^0\;49'$&$+2^0\;2'$\\
$2$&Turuchansk&$65\;55$&$87\;33$&$-9\;19$&$-15\;0$&$+5\;41$\\
$3$&Drontheim&$63\;26$&$10\;24$&$+20\;17$&$+20\;0$&$+0\;17$\\
$4$&Viluisk&$62\;49$&$119\;27$&$+0\;37$&$+1\;52$&$-1\;15$\\
$5$&Bogoslowskoie&$59\;45$&$60\;7$&$-5\;38$&$-9\;9$&$+3\;31$\\
$6$&Fredriksvarn&$59\;0$&$10\;4$&$+20\;18$&&\\
$7$&Jeniseisk&$58\;27$&$92\;11$&$-6\;33$&$-6\;57$&$+0\;24$\\
$8$&Kodiack&$57\;20$&$207\;9$&$-24\;38$&$-26\;43$&$+2\;5$\\
$9$&Copenhagen&$55\;41$&$12\;34$&$+18\;37$&$+17\;40$&$+0\;57$\\
$10$&Altona&$53\;33$&$9\;56$&$+20\;28$&$+18\;43$&$+1\;45$\\[2pt]
$11$&Semipalatinsk&$50\;24$&$80\;21$&$-6\;50$&$-6\;43$&$-0\;7$\\
$12$&Kremsmünster&$48\;3$&$14\;8$&$+18\;26$&$+15\;46$&$+2\;40$\\
$13$&Baker’s Bay&$46\;17$&$235\;58$&$-20\;46$&$-19\;11$&$-1\;35$\\
$14$&Fort Vancouver&$45\;37$&$237\;24$&$-20\;8$&$-19\;22$&$-0\;46$\\
$15$&Toulon&$43\;6$&$5\;55$&$+22\;26$&$+19\;6$&$+3\;20$\\
$16$&Barcelona&$41\;25$&$2\;15$&$+23\;45$&&\\
$17$&Lissabon&$38\;43$&$350\;58$&$+26\;1$&&\\
$18$&Angra (Terceira)&$38\;39$&$332\;47$&$+25\;17$&$+24\;2$&$+1\;15$\\
$19$&Port Bodega&$38\;18$&$236\;58$&$-16\;41$&$-15\;20$&$-1\;21$\\
$20$&Messina&$38\;11$&$15\;34$&$+19\;16$&&\\[2pt]
$21$&Palermo&$38\;7$&$13\;21$&$+19\;29$&$+16\;3$&$+3\;26$\\
$22$&Algier&$36\;47$&$3\;4$&$+23\;18$&$+19\;25$&$+3\;53$\\
$23$&Monterey&$36\;36$&$238\;7$&$-15\;47$&$-14\;13$&$-1\;34$\\
$24$&Gibraltar&$36\;7$&$354\;41$&$+24\;54$&$+21\;40$&$+3\;14$\\
$25$&Zafarine (Ins.)&$35\;11$&$357\;34$&$+24\;35$&$+21\;7$&$+3.\;28$\\
$26$&Sta Barbara&$34\;24$&$240\;19$&$-14\;40$&$-13\;28$&$-1\;12$\\
$27$&San Pedro&$33\;43$&$241\;45$&$-14\;13$&$-13\;8$&$-1\;5$\\
$28$&San Diego&$32\;41$&$242\;47$&$-13\;42$&$-12\;21$&$-1\;21$\\
$29$&San Quentin&$30\;22$&$244\;2$&$-12\;53$&$-12\;6$&$-0\;47$\\
$30$&San Bartolomeo&$27\;40$&$245\;7$&$-12\;1$&$-10\;46$&$-1\;15$\\[2pt]
$31$&Magdalena Bay&$24\;38$&$247\;53$&$-11\;5$&$-9\;15$&$-1\;50$\\
$32$&Mazatlan&$23\;11$&$253\;36$&$-10\;15$&$-9\;24$&$-0\;51$\\
$33$&San Lucas Bay&$22\;52$&$250\;7$&$-10\;31$&$-8\;37$&$-1\;54$\\
$34$&San Blas&$21\;32$&$254\;44$&$-9\;55$&$-9\;0$&$-0\;55$\\
$35$&Socorro Insel&$18\;43$&$249\;6$&$-9\;55$&&\\
$36$&Clarion Insel&$18\;21$&$245\;19$&$-10\;0$&&\\
$37$&Acapulco&$16\;50$&$260\;5$&$-9\;3$&$-8\;23$&$-0\;40$\\
$38$&Trevandrum&$8\;31$&$77\;0$&$-3\;14$&$-0\;44$&$-2\;30$\\
$39$&Cocos Insel&$+5\;53$&$272\;58$&$-8\;11$&$-8\;24$&$+0\;13$\\
$40$&Puna Insel&$-2\;47$&$280\;5$&$-8\;23$&$-8\;56$&$+0\;33$\\[2pt]
$41$&Martins Insel&$-8\;56$&$220\;20$&$-5\;27$&&\\
$42$&Bow Insel&$-18\;5$&$219\;7$&$-5\;21$&&\\
$43$&Rio Grande&$-32\;2$&$307\;40$&$-7\;29$&&\\
$44$&&$-67\;4$&$147\;30$&$+6\;20$&$-12\;35$&$+18\;55$\\
$14$&Sitka&$+57\;3$&$224\;35$&$-28\;45$&$-29\;32$&$+0\;47$\\
$62$&San Francisco&$+37\;49$&$237\;35$&$-16\;22$&$-15\;20$&$-1\;2$\\
$62^{*}$&Oahu&$+21\;17$&$202\;0$&$-12\;19$&$-10\;40$&$-1\;39$\\
$72$&Otaheite&$-17\;29$&$210\;30$&$-5\;45$&$-6\;30$&$+0\;45$\\
\hline
\end{tabular}
\end{center}
\endgroup
\clearpage
\SourceLeaf{185}
\begingroup
\fontsize{9}{10.2}\selectfont
\setlength{\tabcolsep}{1.5pt}
\renewcommand{\arraystretch}{1.02}
\begin{center}
\begin{tabular}{r||r|r|r||r|r|r}
\hline
&\multicolumn{3}{c||}{Inklination}&\multicolumn{3}{c}{Intensität}\\
&Berechn.&Beobacht.&Differenz&Berechn.&Beobacht.&Differenz\\\hline
$1$&$+79^0\;27'$&$+81^0\;9'$&$-1^0\;42'$&$1.675$&&\\
$2$&$77\;20$&$77\;46$&$-0\;26$&$1.662$&$1.678$&$-0.016$\\
$3$&$74\;7$&$74\;12$&$-0\;5$&$1.483$&$1.415$&$+0.068$\\
$4$&$75\;44$&$76\;46$&$-1\;2$&$1.675$&$1.765$&$-0.090$\\
$5$&$70\;45$&$71\;36$&$-0\;51$&$1.556$&$1.524$&$+0.032$\\
$6$&$71\;37$&$72\;1$&$-0\;24$&$1.450$&$1.436$&$+0.014$\\
$7$&$72\;33$&$73\;24$&$-0\;51$&$1.647$&$1.674$&$-0.027$\\
$8$&$73\;22$&$72\;43$&$+0\;39$&$1.638$&$1.603$&$+0.035$\\
$9$&$68\;52$&$70\;0$&$-1\;8$&$1.419$&$1.372$&$+0.047$\\
$10$&$68\;9$&$69\;2$&$-0\;53$&$1.405$&$1.357$&$+0.048$\\[2pt]
$11$&$64\;44$&$65\;18$&$-0\;34$&$1.551$&$1.560$&$-0.009$\\
$12$&$63\;8$&$64\;34$&$-1\;26$&$1.348$&$1.339$&$+0.009$\\
$13$&$71\;12$&$69\;27$&$+1\;45$&$1.675$&$1.643$&$+0.032$\\
$14$&$70\;56$&$69\;22$&$+1\;34$&$1.676$&$1.657$&$+0.019$\\
$15$&$61\;15$&$62\;58$&$-1\;43$&$1.320$&&\\
$16$&$61\;12$&$62\;15$&$-1\;3$&$1.324$&$1.288$&$+0.036$\\
$17$&$63\;0$&$61\;58$&$+1\;2$&$1.352$&$1.299$&$+0.053$\\
$18$&$68\;34$&$66\;50$&$+1\;44$&$1.469$&$1.449$&$+0.020$\\
$19$&$64\;28$&$62\;53$&$+1\;35$&$1.588$&$1.563$&$+0.025$\\
$20$&$54\;12$&$56\;10$&$-1\;58$&$1.219$&$1.232$&$-0.013$\\[2pt]
$21$&$53\;54$&$57\;16$&$-3\;22$&$1.242$&$1.274$&$-0.032$\\
$22$&$56\;52$&$57\;43$&$-0\;51$&$1.267$&$1.272$&$-0.005$\\
$23$&$63\;10$&$61\;4$&$+2\;6$&$1.579$&$1.531$&$+0.048$\\
$24$&$59\;35$&$59\;40$&$-0\;5$&$1.307$&$1.297$&$+0.010$\\
$25$&$57\;32$&$58\;34$&$-1\;2$&$1.283$&&\\
$26$&$61\;23$&$58\;54$&$+2\;29$&$1.559$&$1.501$&$+0.058$\\
$27$&$60\;56$&$58\;21$&$+2\;35$&$1.556$&$1.480$&$+0.076$\\
$28$&$60\;7$&$57\;6$&$+3\;1$&$1.547$&$1.482$&$+0.065$\\
$29$&$57\;42$&$54\;30$&$+3\;12$&$1.514$&$1.461$&$+0.053$\\
$30$&$54\;43$&$51\;41$&$+3\;2$&$1.475$&$1.432$&$+0.043$\\[2pt]
$31$&$51\;24$&$46\;34$&$+4\;50$&$1.434$&$1.362$&$+0.072$\\
$32$&$50\;35$&$46\;38$&$+3\;57$&$1.429$&$1.370$&$+0.059$\\
$33$&$49\;26$&$45\;39$&$+3\;47$&$1.411$&$1.359$&$+0.052$\\
$34$&$48\;35$&$44\;33$&$+4\;2$&$1.405$&$1.362$&$+0.043$\\
$35$&$43\;11$&$40\;44$&$+2\;27$&$1.331$&$1.307$&$+0.024$\\
$36$&$41\;50$&$37\;3$&$+4\;47$&$1.310$&$1.222$&$+0.088$\\
$37$&$+42\;50$&$+37\;57$&$+4\;53$&$1.335$&$1.316$&$+0.019$\\
$38$&$-7\;15$&$-2\;50$&$-4\;25$&$1.014$&$1.012$&$+0.002$\\
$39$&$+27\;46$&$+22\;56$&$+4\;50$&$1.172$&$1.125$&$+0.047$\\
$40$&$+13\;23$&$+9\;8$&$+4\;15$&$1.062$&$1.024$&$+0.038$\\[2pt]
$41$&$-12\;44$&$-14\;6$&$+1\;22$&$1.026$&$1.024$&$+0.002$\\
$42$&$-28\;46$&$-30\;16$&$+1\;30$&$1.125$&$1.123$&$+0.002$\\
$43$&$-33\;14$&$-30\;4$&$-3\;10$&$0.997$&$0.967$&$+0.030$\\
$44$&$-85\;59$&$-87\;30$&$+1\;31$&$2.248$&&\\
$14$&$+76\;30$&$+75\;49$&$+0\;41$&$1.697$&$1.704$&$-0.007$\\
$62$&$+64\;14$&$+62\;6$&$+2\;8$&$1.592$&$1.540$&$+0.052$\\
$62^{*}$&$+37\;36$&$+41\;17$&$-3\;41$&$1.125$&$1.134$&$-0.009$\\
$72$&$-27\;26$&$-30\;18$&$+2\;52$&$1.113$&$1.133$&$-0.020$\\
\hline
\end{tabular}
\end{center}
\endgroup
\clearpage
\SourceLeaf{186}
\{Die Beobachtungen in Palermo wurden Ende 1835 von Dr. \textsc{Sartorius von Waltershausen} und Professor \textsc{Listing} angestellt.

Die Bestimmungen in Gibraltar sowie die der Inklination und Intensität in Algier wurden 1840 auf einer Expedition der norwegischen Korvette Ornen von den Kapitänen \textsc{Konow} und \textsc{Valeur} vorgenommen und uns von Herrn Professor \textsc{Hansteen} mitgeteilt. Die Deklination in Algier wurde 1832 bestimmt und der Description nautique des côtes de l’Algérie par \textsc{Bérard} (Paris 1839) entnommen.

Die Beobachtung auf $67^0\;4'$ südlicher Breite wurde 1840 vom amerikanischen Flottenkapitän \textsc{Wilkes} angestellt und in den Blättern für literarische Unterhaltung 1841, Nr. 6, mitgeteilt.

Die Beobachtungen in Kodiack, Baker’s Bay, Fort Vancouver, Bodega, Monterey, Sta Barbara, San Pedro, San Diego, San Quentin, San Bartolomeo, Magdalena Bay, Mazatlan, San Lucas Bay, San Blas, Acapulco, Cocos Insel und Puna Insel wurden von Kapitän \textsc{Belcher} in den Jahren 1837–1840 vorgenommen und von \textsc{Sabine} in einer der Königlichen Gesellschaft in London vorgelegten Abhandlung \textit{Contributions to terrestical Magnetism} veröffentlicht. Auf Socorro, Clarion, Martins und Bow Island wurden die Deklinationen ebenfalls bestimmt, aber in der Abhandlung von \textsc{Sabine} noch nicht mitgeteilt. Um die noch bestehende Unsicherheit über die Intensität auf Otaheite zu beseitigen, nahm \textsc{Belcher} seinen Rückweg über Otaheite und bestimmte die Elemente auf Point Venus durch zahlreiche Beobachtungen. Neben diesem Ort waren auch Sitka, San Francisco und Oahu, wo \textsc{Belcher} neue Beobachtungen angestellt hat, bereits auf Grundlage anderer Beobachtungen in der ersten Vergleichstabelle enthalten.

Die Elemente von Kremsmünster wurden von Herrn Professor \textsc{Koller} bestimmt. Die Beobachtungen in Trevandrum, die Herr \textsc{Caldecott}, der Leiter des dortigen magnetischen Observatoriums, angestellt hat, werden in einer kleinen Broschüre von \textsc{Sabine} angeführt, \textit{Observations made at the magnetic observatories of Toronto, Trevandrum and St. Helena during a remarkable magnetic disturbance on the 25th and 26th Sept.} 1841.

Die Mitteilung der Beobachtungen in Turuchansk, Drontheim, Viluisk. Bogoslowskoie, Fredriksvarn, Jeniseisk, Copenhagen, Altona, Semipalatinsk, Barcelona, Lissabon, Angra, Messina und Rio Grande verdanke ich der Güte von Herrn Professor \textsc{Hansteen}.\}
